\section{Keeping the same field space in each phase}
\label{ex-phase-fields}

A field label may disappear when one passes to a component of a mixture.
For the selected phase to represent the intended fields, one must rule out
additional null labels. The following argument explains the role of the
finite detectors and the uniform moment estimates used for this purpose.

\begin{lemma}[Moment bounds at every character]\label{ex-character-moments}
Let \(A\in\mathcal A\), let \(t\ge0\), and suppose
\[
 \langle\Omega,(A^*A)^n\Omega\rangle\le t^{2n}
 \qquad(n=0,1,2,\ldots).
\]
Then \(|\chi(A)|\le t\) for every \(\chi\in\Xi\).
\end{lemma}
\begin{proof}
The representing measure gives
\(\int|\chi(A)|^{2n}\,d\mu\le t^{2n}\).
If \(|\chi_0(A)|>t\), choose \(p\) strictly between these numbers.
The nonempty open set \(O=\{\chi:|\chi(A)|>p\}\) has positive measure
\(c\) by full support. Hence
\(cp^{2n}\le\int|\chi(A)|^{2n}\,d\mu\le t^{2n}\).
Since \(p>0\), division gives \(c\le(t/p)^{2n}\to0\), a
contradiction. This also covers \(t=0\), when the right side already
vanishes for positive \(n\).
\end{proof}

Here is the elementary two-point calculation to which these moment
estimates are applied. At a fixed scale \(k\), let \(a,b\) denote two
insertions at distinct points and let \(j_a,j_b,\kappa\in\mathbb C\)
be the prescribed comparison one-point values and connected two-point
value. Form the elements
\begin{align*}
 D_a&=M(a)-j_aI,\qquad D_b=M(b)-j_bI,\\
 D_{ab}&=M(ab)-j_aM(b)-j_bM(a)+(j_aj_b-\kappa)I.
\end{align*}
Assume the preceding lemma's moment bounds hold for each of these three
elements with the same number \(\epsilon_k\ge0\). It follows for every
character that
\begin{align}
 |S_\chi(a)-j_a|&\le\epsilon_k,\label{ex-detector-one}\\
 |S_\chi(ab)-(\kappa+j_aj_b)|
 &\le\epsilon_k(1+|j_a|+|j_b|).\label{ex-detector-two}
\end{align}
Indeed, evaluating \(D_{ab}\) and rearranging gives
\[
 S_\chi(ab)-(\kappa+j_aj_b)
 =\chi(D_{ab})+j_a(S_\chi(b)-j_b)+j_b(S_\chi(a)-j_a),
\]
so the triangle inequality proves \eqref{ex-detector-two}.
If \(\epsilon_k\to0\) and the prescribed one-point values converge,
these errors tend to zero in every phase, with a bound independent of
the character. This conclusion requires moments of the full corrected
operators \(D_{ab}\), not only pointwise estimates for the unselected
two-point correlations.

\begin{proposition}[Finite detectors prevent an additional null space]
\label{ex-no-new-null}
Let \(E\) be a complex vector space of raw labels, and let
\(N\subset N_\chi\) be subspaces. Write
\(\pi:E\to E/N=V\) and \(\pi_\chi:E\to E/N_\chi=V_\chi\).
Let \(W\) be a finite-dimensional complex normed space. Suppose there
are linear maps \(J_k:W\to E\), linear tests
\(\ell_{k,i}:V_\chi\to\mathbb C\), \(1\le i\le r\), and an
injective linear map \(F:W\to\mathbb C^r\) such that:
\begin{enumerate}
\item \(\pi J_k\) is surjective for all sufficiently large \(k\);
\item for every vector of a fixed basis of \(W\),
\[
 F_k w:=\bigl(\ell_{k,i}(\pi_\chi J_kw)\bigr)_{i=1}^r
 \longrightarrow Fw.
\]
\end{enumerate}
Then \(N_\chi=N\).
\end{proposition}
\begin{proof}
Convergence on a finite basis implies convergence in operator norm.
To see this, let \(e_1,\ldots,e_s\) be that basis with continuous
coordinate functionals \(e_j^*\). Then
\[
 \|(F_k-F)w\|
 \le\sum_{j=1}^s|e_j^*(w)|\,\|(F_k-F)e_j\|
 \le\|w\|\sum_{j=1}^s\|e_j^*\|\,\|(F_k-F)e_j\|.
\]
The finite sum tends to zero. If \(W\ne\{0\}\), compactness of its
unit sphere and injectivity of \(F\) give
\(c=\min_{\|w\|=1}\|Fw\|>0\). Choose \(k\) late enough that
\(\|F_k-F\|<c/2\) and \(\pi J_k\) is onto. Then
\(\|F_kw\|\ge(c/2)\|w\|\), so \(F_k\) is injective.
If \(W=\{0\}\), surjectivity already gives \(V=0\), and the result
is immediate.

Now take \(e\in N_\chi\). By surjectivity choose \(w\in W\) with
\(\pi J_kw=\pi e\). Thus \(J_kw-e\in N\subset N_\chi\), and since
\(e\in N_\chi\), also \(J_kw\in N_\chi\). All tests vanish on
\(\pi_\chi J_kw=0\); hence \(F_kw=0\). Injectivity forces \(w=0\),
so \(\pi e=\pi J_k0=0\), or \(e\in N\). This proves
\(N_\chi\subset N\); the reverse inclusion was assumed.
\end{proof}

In the application at label degree \(d\), \(W\) is the comparison
field space, \(J_k\) consists of its corrected labels at the chosen
scales, and the tests are one- or two-point correlations with fixed
exterior labels. Estimates~\eqref{ex-detector-one}--\eqref{ex-detector-two},
combined with continuity at the limiting distinct test configurations,
give convergence on a basis. The comparison tests separate \(W\), so
their joint map \(F\) is injective. The first hypothesis of the
proposition is the eventual spanning statement for the corrected labels
in the unselected quotient. The same-label construction of
Section~\ref{sec:actual-rank-purchases}, followed by the exact-kernel
argument in Section~\ref{sec:finite-rank-kernel-transport}, supplies
this spanning statement. Section~\ref{sec:actual-phase-detector-moments}
supplies the corrected-operator moment bounds on those trials, and
Theorem~\ref{thm:ex-lattice-all-grade-free-detector} supplies the
separating free tests. These applications retain the same corrected
lifts and all separated correlations when passing to the smooth table.
Lemma~\ref{ex-phase-original-null} supplies $N\subseteq N_\chi$.
The proposition then proves equality in every grade for every character.
Thus the finite-dimensional argument uses the actual comparison and
moment constructions; finite dimensionality alone does not supply them.
